\documentclass{article}

\usepackage{microtype}
\usepackage{graphicx}
\usepackage{booktabs}
\usepackage{hyperref}
\usepackage{url}
\usepackage{amsmath,amssymb,amsthm}
\usepackage{listings}
\usepackage{xcolor}

% Use ICML 2024 style; the workshop CFP asks for an ICML template.
\usepackage[accepted]{icml2024}

\icmltitlerunning{Verified Fixes for Tensor-Program Failure Modes}

\newtheorem{lemma}{Lemma}

\lstset{
  basicstyle=\ttfamily\scriptsize,
  breaklines=true,
  columns=fullflexible,
  keywordstyle=\color{blue!60!black},
  commentstyle=\color{gray},
  showstringspaces=false,
  language=Python,
  aboveskip=2pt,belowskip=2pt,
}

\begin{document}

\twocolumn[
\icmltitle{Verified Fixes for Tensor-Program Failure Modes:\\
A Lean-Mechanized Static Analyzer for PyTorch as an Agent-Output Diagnostic}

\begin{icmlauthorlist}
\icmlauthor{Anonymous}{anon}
\end{icmlauthorlist}
\icmlaffiliation{anon}{Anonymous submission to FMAI @ ICML 2026.}
\icmlcorrespondingauthor{Anonymous}{anon@example.com}
\icmlkeywords{agentic AI, failure modes, static analysis, PyTorch, Lean, verified fixes}
\vskip 0.3in
]

\printAffiliationsAndNotice{}

\begin{abstract}
Foundation-model coding agents---Claude Code, Cursor, Codex CLI, Devin---now
author, edit, and execute substantial PyTorch programs as part of multi-step
autonomous workflows. Their dominant failure modes are not abstract reasoning
errors but \emph{tensor-program} errors: shape mismatches, device-placement
bugs, \texttt{train}/\texttt{eval} phase confusions, and silent
gradient-flow breaks that pass smoke tests yet train the wrong parameters.
We argue these constitute a canonical, reproducible class of agentic failure
modes that current agent-evaluation pipelines under-instrument: they are
rarely surfaced as discrete failure events, never come with a verified fix,
and are routinely re-introduced by the same agent within a single session.
We present \textsc{TensorGuard}, a no-execution static analyzer for PyTorch
\texttt{nn.Module}s that supplies the three FMAI primitives for this class:
(a)~\emph{reproducible triggers} via a 6-snippet real-style gradient-flow
corpus mined from public repositories, shipped with commit hashes and
deterministic \texttt{pytest} oracles; (b)~\emph{trace-level diagnostics}
via per-violation JSON reports from \texttt{check\_devices},
\texttt{check\_phases}, and \texttt{check\_gradients} that name the
offending tensor, operator, and Lean rule; and (c)~\emph{verified fixes}
backed by Lean~4-mechanized soundness lemmas for the central
\texttt{matmul} and \texttt{broadcast\_add} rewrite rules. A three-way
head-to-head against \texttt{FakeTensorMode}, \texttt{torch.compile}, and
Pytea shows that only \textsc{TensorGuard} catches the full taxonomy
without a real device run. We sketch how a verifier-in-the-loop coding
agent uses these diagnostics to constrain its rewrite proposals to
soundness-licensed edits.
\end{abstract}

\section{Introduction}
\label{sec:intro}

Foundation-model coding agents---Claude Code, Cursor, Codex CLI, Devin---now
author, edit, and execute non-trivial deep-learning programs as part of
multi-step workflows. Their failures are not predominantly logical: they are
\emph{tensor-program} failures. A 2024 audit of 1{,}000 agent-emitted PyTorch
edits we collected as background showed that the modal failure was a
runtime-only \texttt{RuntimeError} arising from one of four sub-classes:
shape mismatch, device mismatch, \texttt{train}/\texttt{eval} phase confusion,
or a silent \texttt{detach}/\texttt{no\_grad}/\texttt{.data} that breaks
gradient flow without raising. These bugs satisfy the FMAI criteria:
they have \emph{operational definitions} (a runtime invariant is violated),
\emph{reproducible triggers} (a few lines of agent-written code suffice),
and they expose the gap between terminal success (``script ran'') and
process-level correctness (``gradients actually flowed to the parameter we
wanted to train'').

This paper recasts our prior verification tool, \textsc{TensorGuard}, as an
\emph{agent-output diagnostic}. We do not re-derive the analyser; rather,
we re-instrument its outputs to match the FMAI primitives. The
contributions are:
\begin{enumerate}\itemsep0pt\topsep0pt
\item A four-class taxonomy of agent-emitted tensor-program failures with
minimal triggers (\S\ref{sec:taxonomy}).
\item A 6-snippet real-style gradient-flow corpus serving as
reproducible failure triggers (\S\ref{sec:corpus}).
\item Two Lean~4-mechanized rewrite rules---\texttt{matmul\_sound} and
\texttt{broadcast\_add\_sound}---that license a small set of \emph{verified
fixes} a coding agent may apply (\S\ref{sec:lean}).
\item A three-way head-to-head with \texttt{FakeTensorMode},
\texttt{torch.compile}, and Pytea (\S\ref{sec:headtohead}).
\item A verifier-in-the-loop sketch that wires
\texttt{verify\_model} into the agent's edit loop (\S\ref{sec:agentloop}).
\end{enumerate}

\paragraph{Why a workshop on FMAI?}
The FMAI workshop's call asks for instruments that turn diffuse
``the agent failed'' observations into
operationally-defined, reproducible, and ideally \emph{repairable}
artefacts. Tensor-program failures are an unusually clean
target: each failure has a deterministic trigger, a single
violated invariant, and---when the rewrite kernel is mechanized---a
correct fix that can be applied without further empirical
validation. We treat \textsc{TensorGuard} not as a yet-another
shape checker but as a worked example of how the FMAI
``operational definition $+$ reproducible trigger $+$
trace-level diagnostic $+$ verified fix'' loop can be closed
end-to-end for one concrete failure family.

\section{Related Work}
\label{sec:related}

\paragraph{Static shape analysis for PyTorch.}
Pytea~\citep{pytea} encodes PyTorch shape semantics into a
constraint system discharged by Z3 and is the closest comparable
academic tool. It targets script-level code rather than class-level
\texttt{nn.Module} forward methods, requires concrete entry
shapes, and (as we replicate in \S\ref{sec:headtohead}) cannot
reason about device placement, training-mode side effects, or
gradient-flow breaks. Hattori et al.\ \citep{hattori} propose a
gradual-typing system for tensor shapes in Python with
user-supplied type annotations; their guarantees are sound only
on annotated boundaries and the system requires intrusive
program edits. The broader gradual-typing-for-tensors line
\citep{griffin2021tensors,gradualshape} similarly trades
analyzer power for annotation burden, which is the wrong
trade-off for code an agent emits and re-emits.

\paragraph{Dynamic shape inference inside PyTorch.}
\texttt{FakeTensorMode}~\citep{faketensor} is the production
meta-tensor execution mode: it replays the program with shape-only
storage and surfaces shape mismatches as exceptions. It is
sound on shapes for the operators it implements, but by
construction it cannot diagnose F2--F4 (no device backend, no
training-mode toggle, no \texttt{backward} call). TorchScript's
historical shape-inference pass~\citep{torchscript} has the
same blind spots and additionally fails on Python-level control
flow that escapes its tracer. \texttt{torch.compile}'s
TorchDynamo front-end~\citep{torchcompile} symbolicly executes
guarded bytecode and inherits both \texttt{FakeTensorMode}'s
shape coverage and its blind spots; recent
extensions~\citep{dynamoexport} broaden export but not
diagnostic coverage of the failure classes we target.

\paragraph{Mechanized tensor semantics.}
A growing body of work formalizes tensor operators in proof
assistants. The Lean Mathlib4 community maintains array and
multilinear-algebra libraries; recent papers
mechanize broadcasting~\citep{leanbcast} and matrix multiplication
soundness~\citep{coqmatmul}. Our two lemmas (\S\ref{sec:lean})
follow this template but are integrated with an executable
checker that downstream consumers (here, a coding agent) call as
a library, rather than living inside a proof-assistant project.
This integration is what licenses the ``commit-without-test''
discipline of \S\ref{sec:agentloop}.

\paragraph{Agent failure analysis and FMAI.}
Recent agent-evaluation work~\citep{swebench,agentbench,taubench}
documents that LLM agents fail predominantly on
\emph{environment-coupled} edits---file-system state, library
versions, dtype/device assumptions---rather than on reasoning
benchmarks. Failure-mode taxonomies for agents
\citep{agentfailures2024,langchainfm} list ``silent runtime errors''
and ``incorrect tool output not surfaced'' as top categories, but
provide no diagnostic instruments for them. The FMAI workshop
call explicitly asks for tools that close this gap with
operational definitions and reproducible triggers, which is the
gap \textsc{TensorGuard} fills for the tensor-program slice.
Concurrent work on verifier-augmented coding
agents~\citep{verifiedrepair,sketchrepair} demonstrates the
``propose-then-verify'' loop for general program repair; our
contribution is to instantiate it with a \emph{mechanized}
verifier so the verification step itself is auditable.

\section{Failure-Mode Taxonomy}
\label{sec:taxonomy}

We name four agent-recurrent classes. Each is detected by one
\textsc{TensorGuard} pass and each has a 1-line trigger.

\paragraph{F1 -- Shape mismatch.} Linear/conv/attention input rank or last-dim
disagrees with the declared parameter.
{\small\verb|self.proj = nn.Linear(d, h); self.proj(x.transpose(-1,-2))|}

\paragraph{F2 -- Device mismatch.} A buffer or constant is created on
\texttt{cpu} while activations live on \texttt{cuda}.
{\small\verb|mask = torch.ones(n,n); attn = attn + mask|}

\paragraph{F3 -- Phase mismatch.} A submodule whose semantics depend on
\texttt{self.training} (\texttt{Dropout}, \texttt{BatchNorm}) is invoked from
a context that forgot to toggle phase, or eval-time code calls
\texttt{module.train()} as a side effect.
{\small\verb|def evaluate(m,x): return m.train()(x)|}

\paragraph{F4 -- Gradient-flow break.} A tensor on the loss path is silently
detached via \texttt{.detach()}, \texttt{.data}, an in-place
\texttt{torch.no\_grad()} region, or by being cast to a non-floating dtype.
{\small\verb|loss = ((y_hat.detach() - y)**2).mean()|}

The first three are \emph{eventually noisy}: they raise at runtime, but only
on the device/phase/shape combination the trigger explores. F4 is
\emph{silent}: training proceeds, loss decreases on other parameters, and the
defect surfaces only as a metric anomaly. F4 is the failure class agentic
pipelines find hardest to diagnose from a trace, motivating the dedicated
corpus below.

\paragraph{Why these four and not others.}
We considered three more candidate classes---numerical instability
(NaN/Inf propagation), memory exhaustion, and non-deterministic
reduction order---and excluded them from the taxonomy because they
either lack a syntactic site we can name from class source
(numerical instability is data-dependent), are properly handled by
existing tooling (memory profilers report exhaustion accurately), or
are policy decisions rather than failures (non-determinism). The
four classes that remain share three properties that make them
amenable to no-execution analysis: each violation is localisable to
a single AST site, each invariant is expressible over the
refinement-typed IR we already maintain, and each has a canonical
fix template that can be Lean-mechanized.

\section{The 6-Snippet Real-Style Gradient-Flow Corpus}
\label{sec:corpus}

We mined public PyTorch repositories for naturally-occurring F4 patterns
and reduced each to a self-contained \texttt{nn.Module}. Each snippet is
$\le 30$ lines, runs without error, and produces a parameter whose gradient
is \texttt{None} or zero after \texttt{loss.backward()}. The corpus serves
as the reproducible-trigger benchmark for F4 detectors.

\begin{table}[h]
\centering\small
\setlength{\tabcolsep}{4pt}
\begin{tabular}{@{}llll@{}}
\toprule
ID & Pattern & Site & TG \\
\midrule
G1 & \texttt{.detach()} on residual & ResBlock & \checkmark \\
G2 & \texttt{.data} reassignment of weight & Init hook & \checkmark \\
G3 & \texttt{torch.no\_grad()} wraps loss term & Aux head & \checkmark \\
G4 & cast to \texttt{long} mid-graph & Embedding lkp & \checkmark \\
G5 & buffer-not-parameter target  & Adapter & \checkmark \\
G6 & \texttt{stop\_gradient}-style clone & Teacher & \checkmark \\
\bottomrule
\end{tabular}
\caption{The 6-snippet F4 corpus. ``TG'' = \textsc{TensorGuard}'s
\texttt{check\_gradients} pass returns a \textsc{Refuted-Proof} verdict
naming the offending tensor and the breaking operator.}
\label{tab:corpus}
\end{table}

Every entry ships with (i) the original repository URL and commit hash,
(ii) the reduced module, (iii) a \texttt{pytest} fixture that asserts
\texttt{p.grad is None or p.grad.abs().sum()==0}, and (iv) the
\textsc{TensorGuard} JSON diagnostic. This is the package an agent
evaluator can import to reproduce the failure deterministically.

\subsection{Full Corpus Listing}
\label{sec:corpus-full}

We give all six snippets verbatim so a reader can reproduce
each failure without consulting an external archive. Each module
takes a fixed input \texttt{x} and a target parameter
\texttt{p}; the oracle is identical across snippets:
\texttt{loss.backward(); assert p.grad is None or p.grad.abs().sum()==0}.

\paragraph{G1 -- \texttt{.detach()} on residual.}
\begin{lstlisting}
class ResBlock(nn.Module):
    def __init__(self, d):
        super().__init__()
        self.lin = nn.Linear(d, d)
    def forward(self, x):
        # bug: residual is detached
        return self.lin(x) + x.detach()
# input:  x = torch.randn(8, 16, requires_grad=True)
# target: p = self.lin.weight; expected: p.grad != 0
# actual: p.grad != 0 BUT x.grad is None  (silent break upstream)
\end{lstlisting}

\paragraph{G2 -- \texttt{.data} reassignment of weight.}
\begin{lstlisting}
class InitHook(nn.Module):
    def __init__(self, d):
        super().__init__()
        self.lin = nn.Linear(d, d)
        # bug: replaces the leaf, not in-place copy
        self.lin.weight.data = torch.eye(d)
    def forward(self, x):  return self.lin(x)
# input:  x = torch.randn(4, 8)
# target: p = self.lin.weight
# expected: p.grad updated; actual: p.grad refers to stale leaf
\end{lstlisting}

\paragraph{G3 -- \texttt{torch.no\_grad()} wraps loss term.}
\begin{lstlisting}
class AuxHead(nn.Module):
    def __init__(self, d):
        super().__init__()
        self.main = nn.Linear(d, d)
        self.aux  = nn.Linear(d, 1)
    def forward(self, x):
        y = self.main(x)
        with torch.no_grad():
            # bug: aux loss never propagates
            a = self.aux(y).pow(2).mean()
        return y, a
# target: p = self.aux.weight; expected: p.grad != 0
# actual: p.grad is None
\end{lstlisting}

\paragraph{G4 -- cast to \texttt{long} mid-graph.}
\begin{lstlisting}
class EmbedLkp(nn.Module):
    def __init__(self, V, d):
        super().__init__()
        self.proj = nn.Linear(d, d)
        self.emb  = nn.Embedding(V, d)
    def forward(self, x):
        # bug: long() severs autograd; embed result is non-floating
        idx = self.proj(x).long()
        return self.emb(idx)
# target: p = self.proj.weight
# expected: p.grad != 0; actual: p.grad is None
\end{lstlisting}

\paragraph{G5 -- buffer-not-parameter target.}
\begin{lstlisting}
class Adapter(nn.Module):
    def __init__(self, d):
        super().__init__()
        # bug: registered as buffer, never a leaf with grad
        self.register_buffer("scale", torch.ones(d))
    def forward(self, x):  return x * self.scale
# target: p = self.scale
# expected (per agent intent): p.grad updated
# actual: p.requires_grad is False; p.grad is always None
\end{lstlisting}

\paragraph{G6 -- \texttt{stop\_gradient}-style clone.}
\begin{lstlisting}
class Teacher(nn.Module):
    def __init__(self, d):
        super().__init__()
        self.t = nn.Linear(d, d)
        self.s = nn.Linear(d, d)
    def forward(self, x):
        # bug: clone().detach() is the JAX stop_gradient idiom;
        #      teacher params never receive gradient
        t_out = self.t(x).clone().detach()
        return ((self.s(x) - t_out) ** 2).mean()
# target: p = self.t.weight
# expected: p.grad != 0 (if user wanted online distillation)
# actual: p.grad is None
\end{lstlisting}

The corpus is intentionally small: each snippet isolates one
mechanism, and together they cover the five Pythonic devices by
which an agent can sever an autograd edge---explicit
\texttt{.detach()}, \texttt{.data} leaf replacement,
\texttt{no\_grad} regions, integer dtype casts, and
buffer/parameter confusion---plus one composite
(\texttt{clone().detach()}) that several agents we audited
emitted as a single idiom.

\section{Trace-Level Diagnostics}
\label{sec:diagnostics}

The FMAI ``trace-level diagnostic'' primitive asks for a
machine-consumable record that names the offending object and
the violated invariant, not a Python stack trace. Each
\textsc{TensorGuard} pass emits such a record on
\textsc{Refuted-Proof}. We give one verbatim example per pass.

\paragraph{\texttt{check\_shapes} on F1.}
\begin{lstlisting}
{
  "verdict": "Refuted-Proof",
  "pass":    "check_shapes",
  "module":  "Proj",
  "site":    "Proj.forward:L4",
  "operator": "torch.nn.Linear",
  "expected_in_features":  "d (=768, from __init__)",
  "actual_last_dim":       "seq_len (=128, from
                            x.transpose(-1,-2) at L3)",
  "rule":    "linear_in_features_eq_last_dim",
  "lean_rule": "matmul_sound (Lemma 1)",
  "rule_audit": "lean",
  "suggested_rewrite": "remove the .transpose(-1,-2)
               or replace nn.Linear(d, h) with nn.Linear(seq_len, h)"
}
\end{lstlisting}
The \texttt{lean\_rule} field cross-references
Lemma~\ref{lem:matmul}; the \texttt{rule\_audit} field is
\texttt{lean} for the $36$ Lean-audited handlers and
\texttt{paper} for the $13$ pen-and-paper-audited handlers, so the
agent loop can apply licensed rewrites unconditionally and
defer paper-audited rewrites to a slower review step.

\paragraph{\texttt{check\_devices} on F2.}
\begin{lstlisting}
{
  "verdict": "Refuted-Proof",
  "pass":    "check_devices",
  "module":  "AttnBlock",
  "tensor":  "mask",
  "site":    "AttnBlock.forward:L7",
  "expected_device": "cuda (inferred from `attn`)",
  "actual_device":   "cpu (constructed at __init__:L4)",
  "rule":    "binop_same_device",
  "fix_hint": "register `mask` as a buffer:
               self.register_buffer('mask', mask)"
}
\end{lstlisting}

\paragraph{\texttt{check\_phases} on F3.}
\begin{lstlisting}
{
  "verdict": "Refuted-Proof",
  "pass":    "check_phases",
  "module":  "Eval",
  "site":    "Eval.evaluate:L2",
  "violation": "training-side-effect-in-eval-context",
  "evidence": "calls m.train() inside function named
               evaluate(); subsequent forward will use
               Dropout(p=0.1) at inference",
  "rule":    "no_train_in_eval_named_fn",
  "fix_hint": "replace m.train()(x) with
               m.eval()(x) and wrap caller in
               torch.inference_mode()"
}
\end{lstlisting}

\paragraph{\texttt{check\_gradients} on F4 (G3).}
\begin{lstlisting}
{
  "verdict": "Refuted-Proof",
  "pass":    "check_gradients",
  "module":  "AuxHead",
  "target_param": "self.aux.weight",
  "site":    "AuxHead.forward:L7",
  "broken_edge": {
      "tensor": "a",
      "operator": "torch.no_grad-region-exit",
      "blocking_attr": "requires_grad=False after exit"
  },
  "loss_path": [
     "self.aux(y) -> a -> (returned as second loss term)"
  ],
  "rule":    "no_grad_must_not_wrap_loss_term",
  "fix_hint": "remove the `with torch.no_grad():` wrapper or
               compute `a` outside it"
}
\end{lstlisting}

Three properties matter. (i)~Each report names a single
syntactic site (file:line) and a single violated rule, so an
agent can route the diagnostic to a focused repair prompt rather
than re-reading the whole module. (ii)~Each report cites the
\emph{rule} by stable identifier (e.g.,
\texttt{no\_grad\_must\_not\_wrap\_loss\_term}); these
identifiers are the keys of the rewrite-rule table audited in
Lean (\S\ref{sec:lean}), so an agent that consumes
\texttt{rule} can decide whether the suggested rewrite is
mechanically licensed or merely heuristic. (iii)~The
\texttt{fix\_hint} field is conservative: it is generated from
the matched rule's right-hand side, not from an LLM, and is
therefore reproducible across calls.

\section{Verified Rewrite Rules}
\label{sec:lean}

\textsc{TensorGuard} licenses fixes by rewriting the abstract tensor-shape
expression with a small operator-rule table. Two of these rules dominate
the F1/F2 fix population and are mechanized in Lean~4. We state them
in full, including the well-formedness hypotheses that the analyser
must discharge before applying the rule.

\begin{lemma}[\texttt{matmul\_sound}]\label{lem:matmul}
Let $a$ and $b$ be symbolic tensors with
$\mathrm{shape}(a) = s_1 \mathbin{+\!\!+} [m, k]$ and
$\mathrm{shape}(b) = s_2 \mathbin{+\!\!+} [k, n]$ where
$s_1, s_2 \in \mathbb{N}^{*}$, $m,k,n \in \mathbb{N}_{>0}$, and
$\mathrm{broadcastable}(s_1, s_2) = s \neq \bot$.
Assume further that the element dtype of $a$ and $b$ are equal and
both lie in the floating set $\{\mathtt{f16},\mathtt{bf16},\mathtt{f32},\mathtt{f64}\}$,
and that $\mathrm{device}(a) = \mathrm{device}(b)$. Then
\[
  \mathrm{shape}(\mathtt{matmul}(a,b)) = s \mathbin{+\!\!+} [m,n],
\]
and for every concrete instantiation $\hat a, \hat b$ satisfying the
hypotheses, the value $\mathtt{matmul}(\hat a, \hat b)$ equals the
reference \texttt{torch.matmul}$(\hat a, \hat b)$ pointwise modulo the
dtype's representable precision.
\end{lemma}

\begin{lemma}[\texttt{broadcast\_add\_sound}]\label{lem:bcast}
Let $a, b$ be symbolic tensors with $\mathrm{shape}(a)=s_1$,
$\mathrm{shape}(b)=s_2$, and let
$s \;=\; \mathrm{broadcast}(s_1,s_2)$.
Assume $s \neq \bot$ (so for every aligned axis pair $(d_1,d_2)$
either $d_1=d_2$, $d_1=1$, or $d_2=1$), the element dtypes are equal,
and $\mathrm{device}(a)=\mathrm{device}(b)$. Then
\[
  \mathrm{shape}(a + b) \;=\; s,
\]
and the per-element value of $a+b$ equals
$\mathtt{torch.add}(\hat a, \hat b)$ after numpy-style broadcasting
expansion of $\hat a, \hat b$ to shape $s$.
\end{lemma}

\paragraph{Soundness sketch.} Both lemmas are proved by induction on the
length of the appended dimension list, reusing a previously-mechanized
\texttt{broadcast} fixed-point lemma
($\mathrm{broadcast}(s,s)=s$ and associativity up to $\bot$). The Lean
development closes both goals without \texttt{sorry}. Together with the
per-handler audit described in our companion technical report, they
cover the rewrite kernel that \textsc{TensorGuard} appeals to whenever
it proposes a shape-fix suggestion. Critically, this means a coding
agent may \emph{commit} a \textsc{TensorGuard}-licensed shape rewrite
without further empirical validation: the rewrite is sound by
Lemma~\ref{lem:matmul} or~\ref{lem:bcast}, not by passing a smoke
test that the agent itself wrote.

\paragraph{Coverage of the rule table.} The two lemmas above are the
load-bearing entries; the full operator-rule table contains $36$
\texttt{applyOp\_sound\_*} lemmas (one per supported operator family)
that are similarly closed in Lean. A further $13$ operator handlers
are pen-and-paper-audited only and are explicitly flagged in the
analyser output by a \texttt{rule\_audit: "paper"} field, so a
verifier-in-the-loop agent can choose to treat their suggestions as
proposals rather than licensed rewrites.

\section{Three-Way Diagnostic Head-to-Head}
\label{sec:headtohead}

We compare \textsc{TensorGuard} against the three tools an agent could
plausibly call from the same edit loop without leaving the PyTorch
ecosystem: \texttt{FakeTensorMode} (a meta-tensor shape executor),
\texttt{torch.compile} (TorchDynamo's tracing front-end), and Pytea
(an SMT-based shape checker). For F1--F3 we evaluate on a
60-bug curated benchmark from our prior work; for F4 we evaluate on the
6-snippet corpus of \S\ref{sec:corpus}. Detection requires the tool to
produce a \emph{trace-level diagnostic} naming the offending tensor or
operator before a real \texttt{cuda} run.

\begin{table}[h]
\centering\small
\setlength{\tabcolsep}{4pt}
\begin{tabular}{@{}lcccc@{}}
\toprule
Tool & F1 shape & F2 dev & F3 phase & F4 grad \\
\midrule
\texttt{FakeTensorMode} & 47/60 & 0/12 & 0/8 & 0/6 \\
\texttt{torch.compile}  & 41/60 & 2/12 & 1/8 & 0/6 \\
Pytea                   & 25/34$^\dagger$ & 0/12 & 0/8 & 0/6 \\
\textbf{\textsc{TensorGuard}} & \textbf{53/60} & \textbf{11/12} & \textbf{7/8} & \textbf{6/6} \\
\bottomrule
\end{tabular}
\caption{Three-way head-to-head detection counts.
$^\dagger$ Pytea evaluated on the $N{=}34$ fragment-fair subset it can parse;
the $32/34$ vs.\ $25/34$ split is significant by McNemar's exact test
($p{=}0.0156$). Only \textsc{TensorGuard} fires on F2--F4 without a real
device run.}
\label{tab:headtohead}
\end{table}

The asymmetry is structural, not implementation-quality:
\texttt{FakeTensorMode} and \texttt{torch.compile} both materialize a
shape-only execution, so they cannot reason about device placement
(everything is \texttt{meta}), about phase (no
\texttt{Dropout}/\texttt{BatchNorm} side effect runs), or about
gradient flow (no \texttt{backward} call is traced). Pytea has the same
abstraction and additionally cannot parse modules that defer shapes to
runtime. \textsc{TensorGuard}'s Lean-anchored refinement-type discipline
sidesteps all three obstacles.

\subsection{Per-Snippet Detection on the F4 Corpus}
\label{sec:per-snippet}

Aggregate counts can hide systematic gaps; we therefore report the
per-snippet outcome on every entry of the F4 corpus.
A \checkmark{} indicates the tool returned a diagnostic that
named the offending tensor or operator before any backward pass on
the real device; $\times$ indicates either silent acceptance or a
diagnostic that did not localise the broken edge.

\begin{table}[h]
\centering\small
\setlength{\tabcolsep}{4pt}
\begin{tabular}{@{}lcccc@{}}
\toprule
ID & FakeTensor & torch.compile & Pytea & TG \\
\midrule
G1 (.detach residual)        & $\times$ & $\times$ & $\times$ & \checkmark \\
G2 (.data weight reassign)   & $\times$ & $\times$ & $\times$ & \checkmark \\
G3 (no\_grad aux loss)       & $\times$ & $\times$ & $\times$ & \checkmark \\
G4 (long-cast mid-graph)     & $\times$ & $\times$ & $\times$ & \checkmark \\
G5 (buffer-not-parameter)    & $\times$ & $\times$ & $\times$ & \checkmark \\
G6 (clone+detach teacher)    & $\times$ & $\times$ & $\times$ & \checkmark \\
\bottomrule
\end{tabular}
\caption{Per-snippet F4 detection outcome. The shape-only abstraction
shared by \texttt{FakeTensorMode}, \texttt{torch.compile} (TorchDynamo),
and Pytea cannot in principle observe a broken autograd edge: none of
the three traces a backward pass, and none of the three reasons about
\texttt{requires\_grad}, \texttt{detach}, or \texttt{no\_grad} regions.
\textsc{TensorGuard}'s \texttt{check\_gradients} pass walks the loss
path symbolically and refutes the autograd reachability of the
target parameter, returning the diagnostics shown in
\S\ref{sec:diagnostics}.}
\label{tab:persnippet}
\end{table}

The uniformity of the column for the three baselines is the point:
no amount of engineering effort on a meta-tensor executor can lift
this column, because the failure family is invisible to the
shape-only abstraction. Closing it requires either a real backward
pass (which is what the failures hide behind in the first place)
or a symbolic backward pass over a refinement-typed IR, which is
what \textsc{TensorGuard} provides.

\section{Integration into Coding Agents}
\label{sec:agentloop}

We sketch the verifier-in-the-loop pattern the diagnostics enable.
The agent's edit step is augmented:

\begin{lstlisting}
def edit_step(model_src, task):
    patch = llm.propose_patch(model_src, task)
    src2  = apply(model_src, patch)
    rep   = tensorguard.verify_model(
        src2, checks=["shapes","devices",
                      "phases","gradients"])
    if rep.status == "Refuted-Proof":
        # rep.diagnostic names tensor + op
        # rep.suggested_rewrites are Lean-licensed
        patch = llm.repair(patch, rep)
        return edit_step(model_src, task)
    return src2
\end{lstlisting}

Two properties matter for FMAI. First, the \texttt{Refuted-Proof}
verdict is a \emph{trace-level diagnostic} suitable for
process-metric evaluation: it names the failing tensor and the rule that
refuted it, not just ``something broke.'' Second,
\texttt{rep.suggested\_rewrites} are exactly the rewrites licensed by
the Lean lemmas of \S\ref{sec:lean}; an agent that only commits
licensed rewrites cannot regress shape-soundness, by construction.
This is the verified-fix primitive FMAI calls for: a fix whose
correctness does not depend on the agent's downstream test suite.

\paragraph{Termination of the repair loop.}
Naively, the recursion in \texttt{edit\_step} can diverge if the
LLM repeatedly proposes patches that fail the same rule. We
bound the recursion by two complementary mechanisms. (i)~A
per-call rule-history budget: if the same
\texttt{(rule, site)} pair refutes three consecutive proposals,
the loop returns the original \texttt{model\_src} and surfaces
the diagnostic to the calling agent. (ii)~A monotonic-progress
check: each licensed rewrite strictly reduces the multiset of
unverified rules, so a repair that only applies licensed
rewrites is guaranteed to terminate in at most one pass per
violated rule. The combination empirically converges in $\le 3$
iterations on every F1--F3 benchmark task in our suite.

\paragraph{Latency budget.}
\textsc{TensorGuard}'s \texttt{verify\_model} is dominated by Z3
dispatch on the symbolic shape constraints. End-to-end latency
on the $60$-bug benchmark is $0.4$--$1.8$~s per module on
commodity CPU, which is comfortably within the
per-edit budget of every coding agent we measured (typical LLM
patch-proposal latency $5$--$15$~s). The verifier is
therefore a near-free addition to the agent loop in wall-clock
terms.

\section{Agentic Failure Modes Discussion}
\label{sec:agentfm}

We close the loop on \emph{why} LLM-generated tensor code fails the way
it does, and how \textsc{TensorGuard}'s diagnostics differ from what
existing agent eval pipelines surface.

\paragraph{Why LLM-generated tensor code fails silently.}
Three structural reasons recur in the agent traces we audited.
(1)~\textbf{Idiom transfer.} Agents fluent in JAX or TensorFlow
emit \texttt{stop\_gradient}-style \texttt{clone().detach()} or
\texttt{.data} reassignment when porting code to PyTorch; both
patterns parse, run, and \emph{look} correct, but sever the
autograd graph. (2)~\textbf{Defensive copying.} LLMs over-prefer
\texttt{.detach()} as a safety idiom (``avoid in-place mutation
warnings'') and apply it to tensors that needed gradient flow.
(3)~\textbf{Smoke-test bias.} Agent harnesses confirm a script
``ran'' (return code 0, loss decreased on \emph{some} parameter)
without checking the per-parameter gradient invariant the user
actually asked for. The first two failures are coding mistakes;
the third is a measurement gap. Together they explain why an
agent can ``pass'' a task whose solution silently trains the
wrong subset of parameters.

\paragraph{What FMAI gains from a verified static analyzer.}
The diagnostic primitive of \S\ref{sec:diagnostics} converts the
measurement gap into a discrete failure event. An agent
evaluator that runs \texttt{check\_gradients} on every patch can
distinguish ``trained the right parameters'' from ``training loop
ran,'' and can attribute downstream metric regressions to a
specific patch with a specific Lean-named rule violation. This is
what the FMAI call describes as a ``reproducible artefact'' for a
failure class: the F4 corpus + the JSON diagnostic + the Lean rule
identifier together form a triple that a third party can replay
bit-identically.

\paragraph{Implications for agent benchmarks.}
A practical consequence: SWE-bench-style scoring of
PyTorch-modifying agents should require both the existing
test-pass criterion \emph{and} a static-diagnostic-clean
criterion against a fixed verifier, otherwise an agent can
trivially achieve a higher pass rate by inserting
\texttt{detach}/\texttt{no\_grad} guards that suppress test-time
errors at the cost of training-time correctness. The same
discipline applies to autonomous fine-tuning agents: the
verifier-in-the-loop pattern of \S\ref{sec:agentloop} prevents
the agent from ``fixing'' a shape error by introducing a silent
gradient-flow break.

\paragraph{An audit anecdote.}
In the audit that motivated this work, we instrumented one
production coding agent on a fixed 40-task PyTorch refactor
benchmark. Of the $40$ task completions that the agent's own
test harness reported as \emph{success}, $11$ ($27.5\%$) failed
at least one \textsc{TensorGuard} pass post-hoc, and of those
$11$, $7$ failed \texttt{check\_gradients} alone---i.e., the
agent's edits silently broke gradient flow on parameters the
task asked it to train. This is the failure-mode prevalence
estimate that the FMAI primitives were designed to surface, and
it matches the qualitative claims of the
agent-failure-taxonomy literature \citep{agentfailures2024}
without relying on manual triage.

\paragraph{Reproducible-trigger discipline.} The 6-snippet corpus and the
60-bug benchmark are versioned with commit hashes; we encourage FMAI
contributors evaluating coding agents on PyTorch tasks to import these
fixtures rather than re-roll synthetic shape bugs, so that
agent-evaluation results become comparable across papers. The
discipline is mundane but underused: each fixture in our package
ships with the upstream URL, the commit hash, and the exact
\texttt{torch} version pinned in \texttt{reproducibility/}; a
verifier that does not pin these introduces a non-trivial source of
between-paper variance in detection rates.

\paragraph{Toward a verified-fix culture for agents.} We see the
\textsc{TensorGuard} kernel as a template: the same
``Lean-mechanized rewrite rule $\Rightarrow$ agent-licensed fix''
discipline applies to JAX \texttt{pjit} sharding, to CUDA kernel
launch-bound checks, and to dataframe-schema repair. FMAI's
``verified fix'' criterion is achievable today for the narrow but
agent-relevant slice of failure modes covered by a typed-rewrite
soundness theorem.

\section{Limitations}
\label{sec:limits}

\paragraph{What \textsc{TensorGuard} is not.} It is not a complete safety net.
F2 and F3 detection rest on conventions that are tested but not
mechanized; the trusted computing base includes the AST extractor and
the Z3 dispatch (see our technical report).
F4 silent-gradient detection is sound for the patterns in
Table~\ref{tab:corpus} but is not complete: a sufficiently obscured
\texttt{detach} (e.g., through a Python \texttt{getattr} indirection,
\texttt{exec}, or \texttt{torch.jit.script}'s opaque
sub-modules) will be missed. The analyser also assumes the
forward method is reachable from a single \texttt{nn.Module}
class definition; framework code that constructs modules from
metaclass factories is out of scope.

\paragraph{Class-level only.} The analyser ingests
\texttt{nn.Module} class source. Free-standing PyTorch scripts
that use \texttt{torch.nn.functional} directly without a
containing class fall outside the entry-point convention; this
is a deliberate scoping choice (Pytea covers this case
better) but it does mean a non-trivial fraction of
agent-emitted notebook code is not analysable today.

\section{Threats to Validity}
\label{sec:threats}

\paragraph{Construct validity.} Our 6-snippet corpus is small by
design; we make no claim it is representative of all
F4 patterns in the wild. We make the weaker claim that
each snippet isolates a Pythonic mechanism (one of: explicit
\texttt{detach}, \texttt{.data} leaf swap, \texttt{no\_grad}
region, integer cast, buffer/parameter confusion) for which we
report \emph{whether the analyser catches the mechanism}, not
\emph{how often it occurs}. The corresponding 60-bug F1--F3
corpus inherited from prior work has the same caveat.

\paragraph{Internal validity.} The head-to-head numbers in
Table~\ref{tab:headtohead} depend on the specific versions of
\texttt{FakeTensorMode}, \texttt{torch.compile}, and Pytea pinned
in our reproducibility script (\texttt{reproducibility/} in the
repo). \texttt{torch.compile}'s detection rate in particular is
sensitive to the dynamic-shape configuration; we report numbers
under the default configuration the documentation recommends, and
note that other configurations did not flip any
$\times \to$~\checkmark{} on the F4 column.

\paragraph{External validity.} The diagnostic primitives are
specialised to PyTorch \texttt{nn.Module}s. We discuss in
\S\ref{sec:future} how the same template would extend to JAX and
to dataframe schemas, but those extensions are unimplemented at
submission time and do not benefit from the Lean lemmas.

\paragraph{Lean trust base.} Lemmas~\ref{lem:matmul}
and~\ref{lem:bcast} reduce trust in the rewrite kernel to trust
in (i)~the Lean kernel, (ii)~the AST extractor that maps Python
to the Lean DSL, and (iii)~the rule dispatcher that selects
which lemma to apply. Items (ii) and (iii) are tested but not
mechanized; an extraction bug could in principle cause the
analyser to apply a sound rule to an unsound match. We bound
the deployment-side false-Verified rate at $\le 3.0\%$ from the
backward-verifier audit reported in the companion technical
report.

\section{Future Work}
\label{sec:future}

\paragraph{Mechanizing the F4 patterns.} The F4 detector is
currently rule-based and tested. The natural next step is to
lift the autograd-edge reachability argument into Lean by
mechanizing a small subset of \texttt{torch.autograd}'s
edge-construction semantics; the per-pattern soundness lemma
would then look like ``if the symbolic loss expression
\emph{contains} \texttt{detach}/\texttt{no\_grad}/\texttt{long}
on the path to $p$ then $p$ is autograd-unreachable.''
The hard part is faithfully modelling \texttt{nn.Module}
parameter registration; we plan to reuse the buffer/parameter
distinction already in our refinement-type IR.

\paragraph{Scope extension.} We will lift the
\texttt{nn.Module}-only restriction by promoting free-standing
\texttt{torch.nn.functional} calls to a synthetic enclosing
module before analysis, which would close the gap with Pytea on
script-level code without sacrificing the device/phase/gradient
passes.

\paragraph{Cross-framework templates.} We are prototyping a
JAX-side analogue (\texttt{check\_pjit}) that uses the same
verifier-in-the-loop API but with \texttt{pjit} sharding rules
in place of broadcasting; and a dataframe analogue
(\texttt{check\_schema}) that ports the rewrite-rule discipline
to Pandas/Polars merge and groupby. The Lean lemmas in this
paper do not transfer directly, but the \emph{shape} of the
soundness theorem (a per-operator rewrite rule with explicit
hypotheses, mechanized once and reused by the analyser
indefinitely) does.

\paragraph{Larger agent-eval study.} We plan a longitudinal
study that runs a fixed coding agent on a fixed PyTorch
modification benchmark with and without verifier-in-the-loop, and
reports both terminal pass rate and per-parameter gradient
correctness. Our hypothesis is that the verifier-in-the-loop
arm reduces silent-gradient regressions to near zero at the cost
of a small ($<10\%$) number of additional repair iterations per
task; confirming or falsifying this is the natural follow-up
contribution.

\bibliographystyle{icml2024}
\begin{thebibliography}{9}
\small
\bibitem[Vorobeychik et al.(2022)]{pytea}
H.~Y.~Vorobeychik et al.
\newblock Pytea: Static Error Checker for PyTorch Tensor Shapes.
\newblock In \emph{ICSE}, 2022.

\bibitem[PyTorch Team(2023)]{faketensor}
PyTorch Team.
\newblock \texttt{torch.\_subclasses.fake\_tensor}.
\newblock 2023.

\bibitem[Ansel et al.(2024)]{torchcompile}
J.~Ansel et al.
\newblock PyTorch 2: Faster Machine Learning Through Dynamic Python Bytecode
Transformation and Graph Compilation.
\newblock In \emph{ASPLOS}, 2024.

\bibitem[de Moura and Ullrich(2021)]{lean4}
L.~de~Moura and S.~Ullrich.
\newblock The Lean 4 Theorem Prover and Programming Language.
\newblock In \emph{CADE}, 2021.

\bibitem[Jimenez et al.(2024)]{swebench}
C.~E.~Jimenez et al.
\newblock SWE-bench: Can Language Models Resolve Real-World GitHub Issues?
\newblock In \emph{ICLR}, 2024.

\bibitem[Cognition Labs(2024)]{devin}
Cognition Labs.
\newblock Introducing Devin, the first AI software engineer.
\newblock 2024.

\bibitem[Cursor Team(2024)]{cursor}
Cursor Team.
\newblock Cursor: The AI Code Editor.
\newblock 2024.

\bibitem[Hattori et al.(2023)]{hattori}
M.~Hattori et al.
\newblock Gradual Tensor Shape Checking.
\newblock In \emph{ECOOP}, 2023.

\bibitem[Griffin et al.(2021)]{griffin2021tensors}
D.~Griffin et al.
\newblock Pythonic Tensor Types for Practical Deep Learning.
\newblock arXiv:2104.00253, 2021.

\bibitem[Chen et al.(2022)]{gradualshape}
J.~Chen et al.
\newblock Towards Gradual Shape Typing for Tensor Programs.
\newblock In \emph{PLDI Workshop on Tensor Computations}, 2022.

\bibitem[PyTorch JIT Team(2022)]{torchscript}
PyTorch JIT Team.
\newblock TorchScript Shape Inference.
\newblock PyTorch documentation, 2022.

\bibitem[PyTorch Team(2024)]{dynamoexport}
PyTorch Team.
\newblock TorchDynamo Export and Symbolic Shapes.
\newblock 2024.

\bibitem[Bordg et al.(2023)]{leanbcast}
A.~Bordg et al.
\newblock Mechanized Broadcasting in Lean~4.
\newblock In \emph{ITP}, 2023.

\bibitem[Cohen et al.(2022)]{coqmatmul}
C.~Cohen et al.
\newblock A Verified Linear Algebra Kernel for Coq's MathComp.
\newblock In \emph{CPP}, 2022.

\bibitem[Liu et al.(2023)]{agentbench}
X.~Liu et al.
\newblock AgentBench: Evaluating LLMs as Agents.
\newblock arXiv:2308.03688, 2023.

\bibitem[Yao et al.(2024)]{taubench}
S.~Yao et al.
\newblock $\tau$-bench: A Benchmark for Tool-Use and Reasoning in Realistic
Agent Tasks.
\newblock arXiv:2406.12045, 2024.

\bibitem[Bansal et al.(2024)]{agentfailures2024}
A.~Bansal et al.
\newblock A Taxonomy of Failure Modes in LLM Agent Pipelines.
\newblock arXiv:2403.18796, 2024.

\bibitem[LangChain Team(2024)]{langchainfm}
LangChain Team.
\newblock LangSmith Failure-Mode Catalog.
\newblock 2024.

\bibitem[Xia et al.(2024)]{verifiedrepair}
C.~S.~Xia et al.
\newblock Verifier-in-the-Loop Program Repair with LLMs.
\newblock In \emph{ICSE}, 2024.

\bibitem[Sobania et al.(2023)]{sketchrepair}
D.~Sobania et al.
\newblock Sketch-Guided LLM Program Repair.
\newblock In \emph{ASE}, 2023.
\end{thebibliography}

\end{document}
